% ============================================
% Proof of Theorem 1: WAIT Algorithm Optimality
% ============================================

\section{Proof of Theorem~\ref{thm:wait_heavy_traffic}}\label{appexidx::proof of wait_heavy_traffic}

\paragraph{Proof Outline.}
This proof establishes asymptotic optimality of the WAIT algorithm when output lengths are known. The structure proceeds in two parts:
\begin{enumerate}
    \item \textbf{Single-type case}: We construct a coupled dominating process using the Lindley recursion and apply Kingman's bound for queue length expectations.
    \item \textbf{Multiple-type case}: We extend to $m$ types by applying the coupling construction type-by-type, leveraging that the fixed threshold $n_j$ decouples inter-type dynamics.
\end{enumerate}
Key techniques include: Lindley recursion for queue length bounds \citep[Proposition 6.3]{asmussen2003applied}, coupling construction to dominate the original process, Kingman's bound for negative drift processes \citep{kingman1962queues}, and standard random walk results \citep{strait1974maximum}.

Throughout this appendix, we use $b$ to denote the \emph{batch index} (the count of completed batch iterations), which is distinct from the stage index $s \in \{0, 1, \ldots, l_j'\}$ used in the main text for decode stages. We use $t$ for continuous time and $T$ (or $B$) for the total number of batches over the time horizon.

We first consider the single-type case, then extend to multiple types.

\subsubsection{Single-Type Case}

Consider a single prompt type with threshold $n$. For general scheduling algorithms, analyzing this system in continuous time is challenging: KV cache grows dynamically during decode, arrivals are stochastic, and eviction risks complicate the dynamics.

The key insight of the WAIT policy is that its \emph{threshold mechanism achieves dimensionality reduction}. By waiting until exactly $n$ prompts accumulate before initiating a batch, the algorithm enforces a fixed batch size, which ensures constant processing time $\Delta T$ per batch. This decouples the complex continuous-time dynamics into a tractable discrete-time process. This reduction---from a complex memory-constrained queueing system to a simple random walk analysis---is enabled by the algorithm design, not an inherent simplicity of the problem.

Specifically, let $\tilde{\lambda}$ denote the continuous-time arrival rate. The WAIT policy's fixed batch size implies that arrivals during each batch follow $\text{Poisson}(\tilde{\lambda} \cdot \Delta T)$. Defining $\lambda = \tilde{\lambda} \cdot \Delta T$ (expected arrivals per batch), we can analyze the system as an \emph{embedded Markov chain} observed at batch completion times, indexed by $b \in \mathbb{N}$. Setting $n = \lambda$ corresponds to critical loading. The throughput loss arises from ``stuck'' iterations where insufficient prompts are available.

\begin{lemma}[Queue Length and Stuck Time]\label{lemma::throughput gap, single}
This lemma connects queue accumulation to throughput loss through stuck iterations. Let $X^b$ denote the number of arrivals during batch $b$, where $X^b \sim \text{Poisson}(\lambda)$. Let $W^b$ denote the queue length after batch $b$, with $W^0=0$. The state transition is:
$$W^{b+1}=W^b+X^b-\lambda\cdot \mathbf{1}\{W^b+X^b\geq \lambda\}.$$
Define $B_{\emph{stuck}}$ as the number of stuck iterations (where $W^b+X^b<\lambda$). Then:
$$\mathbb{E}[W^B]= \lambda\cdot\mathbb{E}[B_{\emph{stuck}}].$$
\end{lemma}
\noindent The proof is in Appendix~\ref{appendix_lemma::proof of throughput gap, single}. This lemma connects queue accumulation to throughput loss: stuck iterations directly translate to waiting prompts. To bound $\mathbb{E}[W^B]$, we construct a coupled dominating process that admits tractable analysis via classical queueing techniques.

\begin{lemma}[Coupled Dominating Process]\label{lem:coupled_process, single}
The coupled process starts with a safety buffer ($2\lambda$ vs 0) and maintains dominance throughout. Define $\{\tilde{W}^b\}$ by:
$$
\tilde{W}^0 = 2\lambda, \quad \tilde{W}^{b+1} = \max\{ 2\lambda, \tilde{W}^b + X^b - \lambda \}.
$$
Then $\tilde{W}^b \geq W^b + \lambda$ for all $b\in \mathbb{N}$.
\end{lemma}
\noindent The proof is in Appendix~\ref{appendix_lemma::proof of coupled_process, single}.

\paragraph{Lindley recursion.} To analyze the coupled process, we apply the Lindley recursion representation \citep[Proposition 6.3]{asmussen2003applied}:
\begin{lemma}[Lindley Representation]\label{lem:lindley_rep}
Let $S^k=\sum_{r=1}^k (X^r-\lambda)$ be the partial sum process. Then:
$$\tilde{W}^B=2\lambda+\max(S^B, S^B-S^1,\ldots,S^B-S^{B-1},0).$$
\end{lemma}
\noindent This classical queueing representation is made applicable by the WAIT policy's threshold design. Define $\xi^k = X^k-\lambda$. We use time-reversal symmetry to simplify the maximum expression:
$$\max(S^B, S^B-S^1,\ldots,S^B-S^{B-1},0)=\max_{1\leq k\leq B}\Big(\sum_{r=k}^{B}\xi^r\Big)^+\overset{d}{=}\max_{1\leq k\leq B}\Big(\sum_{r=1}^{k}\eta^r\Big)^+, \quad \eta^r\overset{d}{=}X^1-\lambda.$$
Thus $\tilde{W}^B \overset{d}{=}2\lambda+\max_{1\leq k\leq B}(\sum_{r=1}^{k}\eta^r)^+$.

Since $\tilde{W}^B \geq W^B + \lambda$, we have $\mathbb{E}[W^B]\leq \mathbb{E}[\tilde{W}^B]-\lambda$. Let $Y^B = \max_{1\leq k\leq B}(\sum_{r=1}^{k}\eta^r)^+$. By \cite[Lemma 2]{strait1974maximum}:
\begin{lemma}[Random Walk Maximum]\label{lem:Y_T}
$$\ex{}{Y^B}=\ex{}{\max_{1\leq k\leq B}(S^k)^+}=\sum_{k=1}^B \frac{1}{k}\ex{}{(S^k)^+}.$$
\end{lemma}
\noindent This standard result \citep{strait1974maximum} decomposes the maximum into a sum over batch indices. For $\mathbb{E}[(S^k)^+]$, by the Cauchy-Schwarz inequality and $\text{Var}(S^k) = \lambda k$, we have $\mathbb{E}[(S^k)^+] \leq \sqrt{\lambda k}$. Hence:
$$\mathbb{E}[Y^B] \leq \sqrt{\lambda} \sum_{k=1}^B k^{-1/2} \sim \sqrt{\lambda B} \quad \text{as } B \to \infty.$$

\begin{lemma}[Throughput Gap Bound]\label{lemma::throughput gap}
$$\mathbb{E}[W^B]-\lambda\leq \mathbb{E}[Y^B]\leq  \sqrt{\lambda} \sum_{k=1}^B \frac{1}{\sqrt{k}} \sim\sqrt{\lambda B},\quad B\to \infty.$$
\end{lemma}
\noindent The throughput gap is thus $\frac{1}{B}\mathbb{E}[W^B] = O(B^{-1/2})$.

\paragraph{Asymptotic Regime.} In the asymptotic regime where the time horizon scales as $\zeta \to \infty$:
$$\frac{1}{B^{\zeta}}\mathbb{E}[W^{B,\zeta}]= \frac{\lambda^{\zeta}}{B^{\zeta}} + \frac{(\lambda B)^{\zeta/2}}{B^{\zeta}} = \left(\frac{\lambda}{B}\right)^{\!\zeta} + \lambda^{\zeta/2} B^{-\zeta/2} \to 0 \quad \text{as } \zeta \to \infty.$$

\subsubsection{Multiple-Type Case}

We now extend to $m$ prompt types. Consider the WAIT policy $\pi = [n_1,\ldots,n_m]$, where $n_j$ is the threshold for type-$j$ prompts.

\paragraph{Policy Class $\Pi$.} Define the processing time for a full batch (containing all $m$ types):
\begin{equation}\label{eq:batch_time_def}
\Delta T_{[1,\ldots,m]} =  d_0 + d_1 \cdot \sum_{j=1}^m n_j(l_j^\prime+1)\left( l_j+\frac{l_j^\prime}{2} \right).
\end{equation}

\paragraph{Approaching Load Balance.}
The constraint $\Delta T_{[1,\ldots,m]} \leq n_j/\lambda_j$ (Equation~\ref{eq:policy_constraints}) embodies the "approaching load balance" principle. For each type $j$, expected arrivals during a full batch equal $\lambda_j \cdot \Delta T_{[1,\ldots,m]}$. The constraint ensures arrivals do not exceed batch size $n_j$, preventing queue accumulation and eviction risk. The equilibrium policy $\pi^*$ achieves equality (arrivals = completions), maximizing throughput while maintaining stability.

The policy parameters must satisfy:
\begin{equation}\label{eq:policy_constraints}
\begin{aligned}
&\Delta T_{[1,\ldots,m]} \leq \frac{n_j}{\lambda_j} \text{ for all } j \in [m], \\
&M^{\pi} = \sum_{j=1}^mn_j(l_j^\prime+1)\left(l_j+\frac{l_j^\prime}{2}\right)\geq  M^{*}.
\end{aligned}
\end{equation}
These constraints ensure that for any $\pi \in \Pi$: when all $m$ types are batched together, arrivals do not exceed completions for each type (i.e., $\text{Arrival}_j\leq \text{Completion}_j$). The equilibrium policy corresponds to $\pi^* = [ \frac{n_1^*}{l_1^\prime+1},\,\cdots,\,\frac{n_m^*}{l_m^\prime+1} ]$ when $M^\pi = M^*$.

\paragraph{Notation.} Let $b$ denote the batch index, $J^b$ the set of types in batch $b$, $t_s(b)$ the start time and $t_e(b)$ the end time of batch $b$. The waiting time $t_s(b+1)-t_e(b)>0$ occurs only when no type meets its threshold.

\paragraph{State Transition.} At batch $b+1$:
$$W_{(j)}^{t_e(b+1)}=W_{(j)}^{t_e(b)}+X_{(j)}^{[t_e(b), t_s(b+1)]}+Y_{(j)}^{[t_s(b+1), t_e(b+1)]}- n_j\cdot \mathbf{1}\{j\in J^{b+1}\},\quad j=1,\ldots,m,$$
where $X_{(j)}^{[t_e(b), t_s(b+1)]} \sim \text{Poisson}(\lambda_j \cdot (t_s(b+1)-t_e(b)))$ and $Y_{(j)}^{[t_s(b+1), t_e(b+1)]} \sim \text{Poisson}(\lambda_j \cdot (t_e(b+1)-t_s(b+1)))$ are independent arrival counts during waiting and processing periods.

\paragraph{Batch Processing Time.} Recall from~\eqref{eq:batch_time_def} that $\Delta T_{[1,\ldots,m]}$ denotes the processing time for a full batch. Under fixed thresholds, this time is constant across batches.

\paragraph{Type Decoupling.}
Fixed thresholds $n_j$ enable type decoupling. Each type $j$ has independent arrivals during batch processing: $X^b_{(j)} \sim \text{Poisson}(\lambda_j \cdot \Delta T_{[1,\ldots,m]})$. Arrivals depend only on the type-specific rate $\lambda_j$ and batch time $\Delta T_{[1,\ldots,m]}$, not on other types' queue states. This independence allows us to analyze each type separately and sum the resulting bounds.

\begin{lemma}[Multiple-Type Coupled Process]\label{lemma::multiple-type coupled dominace}
Define a coupled process $\{\tilde{W}^b_{(j)}\}$ indexed by batch $b$:
$$\tilde{W}_{(j)}^0=2n_j, \quad \tilde{W}_{(j)}^{b+1}= \max\{ 2n_j,\, \tilde{W}_{(j)}^b +X_{(j)}^b-n_j \},\quad X^b_{(j)} \sim \text{Poisson}(\lambda_j\cdot \Delta T_{[1,\ldots,m]}).$$
Then $\tilde{W}^b_{(j)}\geq W^b_{(j)}$ for all $b\geq 0$ and $j=1,\ldots,m$.
\end{lemma}
\noindent The proof is in Appendix~\ref{appendix_lemma::proof of multiple-type coupled dominace}. Define $Y^b_{(j)} = X^b_{(j)} - n_j$. The coupled process becomes:
$$\tilde{W}^0_{(j)} = 2n_j,\quad \tilde{W}^{b+1}_{(j)} = \max\{ 2n_j, \, \tilde{W}^b_{(j)}+Y^b_{(j)} \}.$$

\paragraph{Case 1: Zero Drift ($\Delta T_{[1,\ldots,m]} = n_j/\lambda_j$).} When $\lambda_j\cdot \Delta T_{[1,\ldots,m]} = n_j$, we have $\ex{}{Y^b_{(j)}}=0$ (zero drift). By Lemma~\ref{lemma::throughput gap}:
$$\ex{}{W^B_{(j)}}\leq \ex{}{\tilde{W}^B_{(j)}}\leq \lambda_j + \sqrt{\lambda_j B}.$$

\paragraph{Case 2: Negative Drift ($\Delta T_{[1,\ldots,m]} < n_j/\lambda_j$).} When $\ex{}{Y^b_{(j)}}= \lambda_j \Delta T_{[1,\ldots,m]} - n_j <0$, the coupled process has negative drift. Since
$$\text{Var}(Y^b_{(j)}) = \lambda_j \cdot \Delta T_{[1,\ldots,m]},\quad
\left| \ex{}{Y^b_{(j)}} \right| = n_j - \lambda_j\cdot \Delta T_{[1,\ldots,m]},$$
by Kingman's bound \citep{kingman1962queues}:
$$\ex{}{W^B_{(j)}}\leq \ex{}{\tilde{W}^B_{(j)}}\leq 2n_j+ \frac{\lambda_j \cdot \Delta T_{[1,\ldots,m]}}{2\left(n_j - \lambda_j\cdot \Delta T_{[1,\ldots,m]}\right)}.$$

\paragraph{Equilibrium Case ($C = M^*$).} When $\pi^* = [\frac{n_1^*}{l_1^\prime+1},\,\cdots,\, \frac{n_m^*}{l_m^\prime+1}]$, all types satisfy $\Delta T_{[1,\ldots,m]} = n_j/\lambda_j$ (zero drift). Summing over all types:
$$\frac{1}{B}\sum_{j=1}^m \ex{}{W^B_{(j)}}\leq \frac{1}{B}\sum_{j=1}^m \lambda_j +\sum_{j=1}^m \sqrt{\frac{\lambda_j}{B}} = O(B^{-1/2}).$$

\paragraph{Asymptotic Regime.} As $\zeta \to \infty$:
$$\frac{1}{\zeta B}\sum_{j=1}^m \lambda_j +\sum_{j=1}^m \sqrt{\frac{\lambda_j}{\zeta B}} \to 0.$$
The expected end-to-end latency scales as $O(B^{-1/2})$.

\paragraph{Strictly Negative Drift Case.} When $\Delta T_{[1,\ldots,m]} < n_j/\lambda_j$ for all $j$, the throughput gap is $O(B^{-1})$:
$$\frac{1}{B}\sum_{j=1}^m \ex{}{W^B_{(j)}}\leq \frac{1}{B}  \sum_{j=1}^m \left(2n_j+ \frac{\lambda_j \cdot \Delta T_{[1,\ldots,m]}}{2\left(n_j - \lambda_j\cdot \Delta T_{[1,\ldots,m]}\right)}\right) = O(B^{-1}).$$
